\documentclass[11pt]{article}
\usepackage[margin=1in]{geometry}
\usepackage{hyperref}
\usepackage{enumitem}
\usepackage{booktabs}
\usepackage{xcolor}
\title{ProbOS --- Month 1, Week 4\\DP Profiling, C++/OpenMP, PDSL v0.1, Cross-Discipline Examples}
\author{Nisong Monyimba}
\date{\today}

\begin{document}
\maketitle

\section*{Week 4 goal}
Three parallel threads: performance (honest DP profiling, C++/OpenMP
kernel), language (PDSL compiler v0.1), and breadth (three
cross-discipline examples proving the kernel generalises beyond
batteries). Close the month with a retrospective and Month 2 plan.

\section*{Day-by-day plan}

\subsection*{Monday -- DP optimisation profiling}
\begin{itemize}[leftmargin=*]
  \item Profile candidate optimisations honestly: cache SALib's
    \texttt{\_build\_problem()} at \texttt{\_\_init\_\_} (measured:
    1040x), CLT convergence subset trick (measured: 45x)
  \item Also test \texttt{inv\_RT} precompute -- measure it, and report
    the true result (\textasciitilde{}1x, no real speedup) rather than
    assuming or claiming a win
\end{itemize}

\subsection*{Tuesday -- C++ BatteryCell + OpenMP}
\begin{itemize}[leftmargin=*]
  \item Port \texttt{BatteryModel2Cell} to C++20:
    \texttt{cpp/include/kernel/battery\_cell.hpp}
  \item \texttt{MonteCarloEngineOMP}: OpenMP-parallel particle
    propagation
  \item Benchmark against Python: measured 7x speedup (serial C++ vs
    Python, not yet parallel-vs-Python)
\end{itemize}

\subsection*{Wednesday -- PDSL compiler v0.1}
\begin{itemize}[leftmargin=*]
  \item Grammar (Lark) $\to$ AST $\to$ codegen $\to$ compiler pipeline
  \item \texttt{python/pdsl/}: \texttt{grammar.lark},
    \texttt{ast\_nodes.py}, \texttt{parser.py}, \texttt{codegen.py},
    \texttt{compiler.py}
  \item Literature update: added modern citations (Feng 2018, Saltelli
    2019, Chopin \& Papaspiliopoulos 2020, Naesseth 2019, FDA 2019) to
    the manuscript
\end{itemize}

\subsection*{Thursday -- v0.1.0 checkpoint decision}
\begin{itemize}[leftmargin=*]
  \item Decision: defer formal publication/Zenodo DOI/GitHub Release
    to Year 2 -- a git tag is version-control hygiene, a Zenodo DOI is
    publication infrastructure with zero value before there is a paper
    citing it
  \item Repo hygiene: moved all generated PNGs into
    \texttt{outputs/figures/}, fixed resulting CI break (missing
    \texttt{lark} dependency)
  \item \texttt{check\_ci.sh}: permanent CI-watching script
\end{itemize}

\subsection*{Friday/Saturday -- Three cross-discipline examples}
\begin{itemize}[leftmargin=*]
  \item \texttt{week4\_option\_pricer.py} -- European call option
    (Black-Scholes-Merton GBM), validated against the closed-form price
  \item \texttt{week4\_ed\_queue.py} -- M/M/1 emergency department
    queue, validated against closed-form queueing theory ($\rho$, $L$,
    $W$)
  \item \texttt{week4\_clinical\_trial.py} -- adaptive Bayesian trial
    (Beta-Binomial), informed by FDA (2019) adaptive design guidance
  \item Each needed zero kernel changes -- only a new \texttt{Model}
    subclass, proving domain-agnosticism rather than assuming it
\end{itemize}

\subsection*{Sunday -- Month 1 retrospective + Month 2 plan}
\begin{itemize}[leftmargin=*]
  \item Wrote \texttt{docs/retrospectives/month1\_retrospective.md} and
    \texttt{month2\_plan.md}
  \item Compiled both to PDF under
    \texttt{docs/monthly\_plans/month1/} and
    \texttt{docs/monthly\_plans/month2/}
  \item Created \texttt{docs/monthly\_plans/overall/main.tex}: replaced
    the earlier inconsistent ``Year 1/2/3'' labelling (which did not
    align with actual calendar months) with an honest Month 1--24
    numbering, most of it explicitly DIRECTIONAL rather than
    over-specified
\end{itemize}

\section*{What was actually accomplished (retrospective)}

\begin{itemize}[leftmargin=*]
  \item DP profiling: two genuine wins (1040x, 45x) and one honest
    non-win (\texttt{inv\_RT} \textasciitilde{}1x) -- documented
    rather than glossed over
  \item C++ \texttt{BatteryCell} + \texttt{MonteCarloEngineOMP}: 7x
    serial speedup over Python confirmed by measurement
  \item PDSL compiler v0.1: full grammar $\to$ AST $\to$ codegen $\to$
    compiler pipeline, 28 tests, including the tricky unary-minus-in-
    function-call parser case
  \item Three cross-discipline examples, each validated against a
    closed-form or literature benchmark, not just ``runs without
    crashing''
  \item \texttt{check\_ci.sh} established as a permanent, repeatable
    CI verification step
  \item Month 1 retrospective and Month 2 plan written and compiled to
    PDF, plus the overall 24-month roadmap correcting earlier
    inconsistent year labelling
  \item \emph{Later, before Week 7:} a full pre-week quality audit
    system was added retroactively --
    \texttt{docs/standards/quality\_standards.md} and
    \texttt{pre\_week\_audit.sh} -- covering coverage, property-based
    testing (Hypothesis), security scanning (bandit, pip-audit), and
    packaging/build verification. This audit caught and fixed real
    latent bugs from Week 4's own work: \texttt{python/examples/}
    missing \texttt{\_\_init\_\_.py}, 4 mypy errors that had never been
    checked by CI, a stale coverage threshold, and a bandit finding on
    the PDSL compiler's \texttt{exec()} call (verified as a sound,
    documented security boundary rather than a real vulnerability)
\end{itemize}

\section*{Numbers at Week 4 close (updated through the pre-Week-7 audit)}
\begin{itemize}[leftmargin=*]
  \item Python tests: 322/322 passing
  \item C++ tests: 13/13 passing
  \item mypy strict: 0 errors (full \texttt{python/} tree)
  \item ruff: 0 warnings (full \texttt{python/} tree)
  \item Coverage: 90.79\% (floor: 85\%)
  \item bandit: 0 unresolved findings
  \item pip-audit: 0 known CVEs
  \item \texttt{pre\_week\_audit.sh}: 21/21 checks passed
\end{itemize}

\section*{Carried into Month 2}
Month 1 closes with a validated forward-simulation kernel across four
domains (battery physics, finance, hospital operations, clinical
medicine), a working compiler front-end, and -- critically -- a
standing quality-audit protocol now running automatically in CI before
every future week's work begins. Month 2 turns to sequential inference
(\texttt{ParticleFilter}), binding the C++ kernel into the Python
package (pybind11), and eventually a FastAPI service layer.

\end{document}
